% error-handling.tex — throws/try/try?/try!, do-catch patterns, rethrows,
% Result, typed throws (Swift 6), defer, errors across async/Task,
% cancellation, LocalizedError, never-crash rules.
% Sources (checked 2026-10-06):
%   docs.swift.org/swift-book Error Handling (throws, do-catch, try? / try!, rethrows, defer
%     cannot transfer control out) + Assertions and Preconditions (assert not evaluated in -O;
%     precondition evaluated in -O, not in -Ounchecked)
%   swift-evolution SE-0413 (typed throws, 6.0, do throws(E)) · SE-0304 (structured
%     concurrency: a throwing child cancels its siblings) — checked against the evolution index
%   developer.apple.com/documentation/swift/task/result · /foundation/localizederror
% @labels: area=swift kind=concept platform=apple topic=language discussion=145
% @tags: throws, rethrows, do-catch, try-optional, result-type, typed-throws, defer, cancellationerror, localizederror, withthrowingtaskgroup, precondition
\documentclass[8pt]{extarticle}
\usepackage{printup-sheet}
\usepackage{array}

\lstdefinelanguage{SwiftSheet}{
  morekeywords={protocol,class,final,struct,enum,func,var,let,init,case,
    if,else,return,guard,self,nil,private,true,false,in,async,await,throws,
    rethrows,try,do,catch,throw,defer,where,is,as,switch,Task,some},
  sensitive=true, morecomment=[l]{//}, morestring=[b]"}

\tikzset{
  fr/.style={draw=sheetBlue, thick, fill=sheetBlue!6, rounded corners=2pt,
             minimum width=27mm, minimum height=6mm, font=\ttfamily\scriptsize,
             align=left, anchor=west},
  lbl/.style={font=\scriptsize, text=black!80, inner sep=1pt, align=center},
  err/.style={->, very thick, draw=sheetRed},
  cc/.style={draw=sheetGrey, fill=white, rounded corners=2pt, font=\ttfamily\scriptsize,
             inner sep=2pt, anchor=west},
}
\newcolumntype{L}[1]{>{\raggedright\arraybackslash}p{#1}}

\begin{document}

\sheettitle{Error handling — throws, Result, typed throws, async}{swift · folio}

\oneliner{Swift errors are \textbf{values} (any type conforming to the empty
protocol \texttt{Error}) returned on a \textbf{separate, checked path}: a
function says \texttt{throws}, every call site says \texttt{try}, and the error
travels up until a \texttt{do/catch} handles it. No stack unwinding across
unmarked frames, no cost on the success path, no stack trace.}

\begin{multicols}{2}
\raggedright

\section{How it works}
\begin{itemize}
  \item \textbf{\texttt{try}} propagates (caller must \texttt{throw} or
        \texttt{catch}); \textbf{\texttt{try?}} turns the error into
        \texttt{nil} and \emph{flattens} (\texttt{throws -> Int?} gives
        \texttt{Int?}, since Swift 5); \textbf{\texttt{try!}} traps on error.
  \item \textbf{\texttt{catch} clauses} are patterns, checked \textbf{top to
        bottom}: \texttt{catch E.timeout}, \texttt{catch E.status(let c) where c >= 500},
        \texttt{catch let e as DecodingError}, \texttt{catch is CancellationError};
        a bare \texttt{catch} binds \texttt{error}. In a non-throwing function the
        catches must be \textbf{exhaustive} (usually: end with bare \texttt{catch}).
  \item \textbf{\texttt{rethrows}}: throws only if its closure argument does
        — \texttt{map}, \texttt{filter} need \texttt{try} only for a throwing closure.
  \item \textbf{\texttt{Result<Success, Failure: Error>}} = the outcome as a
        value: store it, pass it to a callback, collect many.
        \texttt{Result \{ try f() \}} $\leftrightarrow$ \texttt{try r.get()};
        \texttt{map}, \texttt{mapError}, \texttt{flatMap}.
  \item \textbf{Typed throws} (Swift 6, SE-0413): \texttt{throws(ParseError)}.
        \texttt{throws} $\equiv$ \texttt{throws(any Error)}; non-throwing
        $\equiv$ \texttt{throws(Never)}. A \texttt{do} whose calls all throw
        \texttt{E} gives \texttt{catch} an \texttt{error: E} — exhaustive
        \texttt{switch}, no casting. Generic \texttt{throws(E)} replaces
        \texttt{rethrows}. Advice: \textbf{keep public APIs untyped} (adding
        a case later breaks clients); typed for closed internal domains/Embedded.
  \item \textbf{\texttt{defer}} runs on \emph{every} scope exit, \textbf{LIFO};
        after a \texttt{return} expression is evaluated; it cannot
        \texttt{throw} or \texttt{return} itself.
  \item \textbf{\texttt{LocalizedError}}: \texttt{errorDescription},
        \texttt{failureReason}, \texttt{recoverySuggestion} — feeds
        \texttt{localizedDescription}. \texttt{CustomNSError}: domain/code for
        \texttt{NSError} bridging.
\end{itemize}

\section{Example}
\begin{lstlisting}[language=SwiftSheet]
enum NetErr: Error { case timeout, status(Int) }
func refresh() async {                 // non-throwing
  isLoading = true
  defer { isLoading = false }          // on every exit
  do {
    let data = try await api.fetch()   // throws NetErr
    items = try JSONDecoder().decode([Item].self, from: data)
  } catch NetErr.status(let c) where c >= 500 {
    banner = "Server down (\(c))"
  } catch is CancellationError {
    return                             // user left: no banner
  } catch {                            // required catch-all
    log(error); banner = error.localizedDescription
  }
}
\end{lstlisting}

\section{Never-crash rules}
\begin{itemize}
  \item No \texttt{try!} / \texttt{!} on input you do not own (network, disk,
        user). \texttt{try!} only for invariants: a bundled fixture, a literal regex.
  \item \texttt{assert} = debug only (gone in \texttt{-O});
        \texttt{precondition} stays in release; \texttt{fatalError} always.
        All three are for \emph{programmer} errors, not runtime conditions.
  \item Wrap with context, never swallow:
        \texttt{catch \{ throw AppError.parse(underlying: error) \}}.
\end{itemize}

\columnbreak

\section{How an error travels}
\begin{tikzpicture}[sheet]
  % frames (top = caller)
  \node[fr] (f1) at (0,0)    {refresh()\\\textcolor{sheetGreen}{do \{ \ldots\ \} catch}};
  \node[fr] (f2) at (0,-1.0) {fetch() throws\\\textcolor{sheetBrown}{defer \{ close() \}}};
  \node[fr, draw=sheetRed, fill=sheetRed!6] (f3) at (0,-2.0) {decode() throws\\\textcolor{sheetRed}{throw E.bad}};
  \draw[err] ([xshift=-3mm]f3.north east) -- node[lbl, right]{(1) \texttt{try}} ([xshift=-3mm]f2.south east);
  \draw[err] ([xshift=-3mm]f2.north east) -- node[lbl, right]{(3) \texttt{try}} ([xshift=-3mm]f1.south east);
  \node[lbl, text=sheetBrown, anchor=west] at (2.85,-1.0) {(2) its \texttt{defer}s run, LIFO};
  \node[note, anchor=west] at (2.85,-2.0) {every frame on the way says \texttt{try}:\\the path is visible in source};
  % catch ladder
  \node[cc] (c1) at (4.1,0.55) {catch E.timeout};
  \node[cc] (c2) at (4.1,0.0)  {catch let e as DecodingError};
  \node[cc] (c3) at (4.1,-0.55) {catch \{ error \}};
  \draw[err] (f1.east) -- (c1.west);
  \draw[flow, sheetGrey] (c1.south west) ++(0.25,0) -- ++(0,-0.21);
  \draw[flow, sheetGrey] (c2.south west) ++(0.25,0) -- ++(0,-0.21);
  \node[lbl, text=sheetGreen, anchor=west] at ([xshift=2pt]c3.east) {top-down: first match wins};
  % Task boundary
  \node[draw=sheetOrange, very thick, dashed, rounded corners=3pt, minimum width=79mm,
        minimum height=10mm, anchor=north west] (tb) at (0,-2.65) {};
  \node[font=\bfseries\scriptsize, text=sheetOrange, anchor=north west] at (tb.north west)
    {Task boundary: \texttt{let t = Task \{ try await load() \}} keeps the error inside};
  \node[lbl, anchor=west] at (0.1,-3.3) {\texttt{try await t.value}: rethrows};
  \node[lbl, anchor=west] at (3.05,-3.3) {\texttt{await t.result}: \texttt{Result}};
  \node[lbl, anchor=west, text=sheetRed] at (5.6,-3.3) {never read: \textbf{lost}};
\end{tikzpicture}

\section{Across async \& Task}
\begin{itemize}
  \item \texttt{async throws} composes: call with \texttt{try await}. Order of
        keywords: \texttt{try await f()}.
  \item Unstructured \texttt{Task\{\}} stores the error in
        \texttt{Task<T, any Error>}; nothing forces you to read it.
  \item \texttt{async let}: error surfaces at the \texttt{await}; siblings are
        cancelled when the scope exits. \texttt{withThrowingTaskGroup}: first
        thrown child error leaves the group $\to$ the rest are cancelled.
  \item \textbf{Cancellation} is cooperative: \texttt{try Task.checkCancellation()}
        and \texttt{Task.sleep} throw \texttt{CancellationError};
        \texttt{URLSession} throws \texttt{URLError(.cancelled)} instead —
        treat both as ``not a failure''.
\end{itemize}

\section{Traps}
\begin{itemize}
  \trap{\texttt{try?} on \texttt{throws -> Int?}: \texttt{nil} means
        ``threw'' \emph{or} ``returned nil'' — use \texttt{do/catch}.}
  \trap{In \texttt{do \{ try a(); try b() \}}, if \texttt{a} throws,
        \texttt{b} \textbf{never runs}.}
  \trap{A \texttt{defer} that mutates the returned variable does
        \textbf{not} change the already-evaluated return value.}
  \trap{Typed throws in a public API = a breaking change for every new case.}
\end{itemize}

\section{Remember}
\textbf{``try marks, throw sends, catch ends; defer cleans backwards.''}
\texttt{try?} = forget, \texttt{try!} = bet the app, \texttt{Result} = keep for later.


\end{multicols}

\noindent{\footnotesize\color{sheetGrey}\textit{Related:} optionals · async/await \& cancellation · continuations (\texttt{resume(throwing:)}) · NSError bridging · Combine \texttt{Failure} · logging \& crash reporting}

\end{document}
